\begin{figure}[htbp]
    \raggedright
    \resizebox{0.96\textwidth}{!}{%
    \begin{tikzpicture}[
        font=\small,
        title/.style={font=\bfseries\large},
        electrode/.style={draw=blue!55!black, fill=blue!15, minimum width=4mm, minimum height=1.8mm},
        coil/.style={draw=purple!55!black, fill=purple!15, minimum width=4mm, minimum height=6mm},
        potential/.style={red!70!black, thick, opacity=.55},
        field/.style={purple!75!black, thick, -{Latex[length=2.3mm]}},
        legendtext/.style={font=\scriptsize, anchor=west, align=left},
        arrow/.style={-{Latex[length=2.5mm]}, thick}
    ]
        % EEG panel
        \node[title] at (-4.6,3.45) {EEG};
        \draw[thick] (-4.6,0.8) circle (2.0cm);
        \draw[fill=gray!12] (-4.6,0.8) circle (1.45cm);
        \draw[fill=red!35, draw=red!60!black] (-4.9,0.55) circle (3mm);
        \node[font=\footnotesize, align=center] at (-4.6,-2.15) {Potenciales eléctricos\\difusos en el cuero cabelludo};
        \foreach \a in {20,55,90,125,160} {
            \node[electrode, rotate=\a-90] at ($(-4.6,0.8)+(\a:2.08cm)$) {};
        }
        \foreach \a in {200,235,270,305,340} {
            \node[electrode, rotate=\a-90] at ($(-4.6,0.8)+(\a:2.08cm)$) {};
        }
        \draw[potential] (-4.9,0.55) -- ($(-4.6,0.8)+(20:2.0cm)$);
        \draw[potential] (-4.9,0.55) -- ($(-4.6,0.8)+(55:2.0cm)$);
        \draw[potential] (-4.9,0.55) -- ($(-4.6,0.8)+(125:2.0cm)$);
        \draw[potential] (-4.9,0.55) -- ($(-4.6,0.8)+(200:2.0cm)$);
        \draw[potential] (-4.9,0.55) -- ($(-4.6,0.8)+(305:2.0cm)$);

        % MEG panel
        \node[title] at (4.6,3.45) {MEG};
        \draw[thick] (4.6,0.8) circle (2.0cm);
        \draw[fill=gray!12] (4.6,0.8) circle (1.45cm);
        \draw[purple!35, thick] (4.6,0.8) circle (2.38cm);
        \draw[fill=red!35, draw=red!60!black] (4.3,0.55) circle (3mm);
        \node[font=\footnotesize, align=center] at (4.6,-2.15) {Campos magnéticos\\medidos por sensores externos};
        \foreach \a in {25,60,130,165,205,245,325} {
            \node[coil, rotate=\a-90] at ($(4.6,0.8)+(\a:2.5cm)$) {};
        }
        \draw[field] (4.3,0.55) .. controls (3.35,1.35) and (3.55,2.35) .. (4.65,2.78);
        \draw[field] (4.3,0.55) .. controls (5.45,1.25) and (5.85,2.15) .. (5.25,2.72);
        \draw[field] (4.3,0.55) .. controls (5.55,0.15) and (5.75,-0.75) .. (4.95,-1.05);
        \draw[field] (4.3,0.55) .. controls (3.25,0.0) and (3.45,-0.95) .. (4.25,-1.16);

        \draw[dashed, thick, gray] (0,-1.4) -- (0,3.25);
        \node[align=center, font=\small] at (0,1) {Misma actividad\\neuronal subyacente};
        \draw[arrow] (-0.6,0.5) -- (-2.1,0.5);
        \draw[arrow] (0.6,0.5) -- (2.1,0.5);

        % Legend
        \node[legendtext, font=\scriptsize\bfseries] at (-5.8,-3.0) {Leyenda};
        \draw[thick] (-5.55,-3.35) circle (1.5mm);
        \node[legendtext] at (-5.25,-3.35) {cabeza/cuero cabelludo};
        \draw[fill=gray!12] (-5.55,-3.75) circle (1.5mm);
        \node[legendtext] at (-5.25,-3.75) {cerebro};
        \draw[fill=red!35, draw=red!60!black] (-5.55,-4.15) circle (1.5mm);
        \node[legendtext] at (-5.25,-4.15) {fuente neuronal};
        \draw[potential] (-5.75,-4.55) -- (-5.35,-4.55);
        \node[legendtext] at (-5.25,-4.55) {proyección eléctrica};

        \node[electrode, minimum width=4mm, minimum height=1.8mm] at (0.7,-3.35) {};
        \node[legendtext] at (1.05,-3.35) {electrodos EEG};
        \node[coil, minimum width=4.8mm, minimum height=3.2mm] at (0.7,-3.75) {};
        \node[legendtext] at (1.05,-3.75) {sensores MEG};
        \draw[purple!35, thick] (0.7,-4.15) circle (1.5mm);
        \node[legendtext] at (1.05,-4.15) {casco/sensores externos};
        \draw[field] (0.4,-4.55) -- (1.0,-4.55);
        \node[legendtext] at (1.05,-4.55) {campo magnético};
    \end{tikzpicture}}
    \caption{Esquema conceptual de la observación de una fuente neuronal mediante EEG y MEG.}
    \label{fig:eeg_meg_concept}
\end{figure}
